\chapter{Results}

For the execution of the project, some development steps had to be followed: familiarization with the system, study of the involved modules, reading of the requirements, elaboration of a document describing the entire implementation process and relationship with the various modules, implementation and tests.


\section{Project Activities}
\label{metodo3}

\section {System Requirements}
\label{req}


\section{Development and Implementation}

The Table \ref{tab:table} presents the executed activities.

\begin{table}
\centering
\begin{tabular}{|c|r|l|}\hline
		Activity 1 & aa  & ab  \\
					 & a & b \\ \hline
		Activity 2  & aa & ab \\
					 &  a & b \\ \hline
		\end{tabular}
	\caption{Example table - Put all information about the table here}
	\label{tab:table}
\end{table}

\section{Tests}

\section{Analysis of the Results}

Present the results without adulterations and make objective analyses. Think about the best way to present the results graphically. If the graphs are difficult to interpret, perhaps tables would be a better way to present results. Do not present data (graphs and tables) if there is not an interesting conclusion to be drawn. Remember to be concise.

\emph{Do not forget the units!} Think that \emph{a priori} every number should have a unit. Do not write the units in italics (in the mathematical environment) and be careful to differentiate uppercase and lowercase letters. An example is to write $22$ [kN] and not $22 KN$ (Kelvin times Newton!).

When presenting experimental results, be careful to also present the calculation of uncertainties whenever they are significant. When making conclusions, always consider whether the size of your sample is large enough from a statistical point of view. Remember that the empirical mean $\hat{\mu}_X$ of $N$ independent observations of the variable $X_i$ has variance
\[
\hat{\sigma}_{\mu}^2 = \frac{1}{N(N-1)} \sum_{i=1}^N (X_i-\hat{\mu}_X)^2
\enspace,
\]
where it is assumed that the variables $X_i$ have the same distribution and that this distribution has a finite second moment.

\section{Chapter Summary}
\label{sec:summary4}
This section should be a summary of the chapter, highlighting the most important points. It should be concise and objective, and it should not contain new information that has not been presented in the chapter.

\clearpage
